\subsection{The concrete two-end bounds for periodic billiards}
\label{sec:v77-determinant-bridge}

For completeness we expose the hypotheses of the separated determinant
comparison at the point where the periodic Jacobi estimates supply them.
This also distinguishes the scalar twist $D_{N,b}$ from a block operator.
Use $\mathsf B_N$ for the latter.  The trace-ideal norms and parameter
norms have the meanings of Proposition~\ref{prop:v76-determinant}.

\begin{corollary}[Periodic end blocks with explicit rates]
\label{cor:v77-periodic-blocks}
Fix a finite differentiation order $k$ in Theorem
\ref{thm:v64-periodic-relative}.  On a sufficiently small common collar
there are $0<q_G<\rho<1$, $q_*<1$, $a_k<\infty$ and $C_k$ as
follows.  Put $r=\lfloor N/3\rfloor$, $N\ge6$, let $P_-,P_+$
project onto the first and last $r$ interior sites, and let $P=P_-+P_+$.
Write $G_N=H_N(0)^{-1}$, $\Delta_N=H_N(y)-H_N(0)$ and
$\Delta_E=P\Delta_NP$.  The determinant of
$I+G_N\Delta_E$ reduces by Sylvester's identity to the two-end matrix
\[
 T_E=(PG_NP)\Delta_E=\mathsf B_N+E_N,
\]
with $\mathsf B_N=\diag(T_{N,-},T_{N,+})$ block diagonal and $E_N$ block off diagonal.
There are trace-class half-line operators $T_-(u),T_+(v)$ such that
\begin{align}
 \|G_N(\Delta_N-\Delta_E)\|_{C^k(\mathfrak S_1)}
       &\le C_k(1+N)^{a_k}\rho^r,\label{eq:v77-middle-rate}\\
 \|\mathsf B_N\|_{C^k(\mathfrak S_1)}&\le C_k,\notag\\
 \|E_N\|_{C^k(\mathfrak S_2)}
       &\le C_k(1+N)^{a_k}q_G^{N-2r},\label{eq:v77-cross-rate}\\
 \|T_{N,\pm}-T_\pm\|_{C^k(\mathfrak S_1)}
       &\le C_k(1+N)^{a_k}
          \{\rho^r+\rho^{N-r}+q_G^{2(N-r)}\}.
       \label{eq:v77-diagonal-rate}
\end{align}
In the last line finite blocks are padded by zero in their half-line
spaces.  The relevant interpolation segments have operator norm at most
$q_*$.  Consequently, with $\ell_\pm=\log\det(I+T_\pm)$,
\begin{equation}\label{eq:v77-concrete-determinant}
 \|\log\det(I+G_N\Delta_N)-\ell_--\ell_+\|_{C^k}
 \le C_k(1+N)^{b_k}
       \{\rho^r+\rho^{N-r}+q_G^{2(N-r)}+q_G^{2(N-2r)}\}.
\end{equation}
The exponents $a_k,b_k$ are finite fixed-order polynomial losses;
no dimension factor is hidden in the trace-ideal comparison.
\end{corollary}
\begin{proof}
The proof of Theorem~\ref{thm:v64-periodic-relative} gives the local
entry bounds on $\Delta_N$ and its fixed derivatives by the endpoint
weights $\rho^i+\rho^{N-i}$, up to fixed polynomial factors after
gluing comparisons.  For a tridiagonal matrix the sum of absolute
entries bounds trace norm.  Summing entries touching the deleted middle
therefore gives $C_k(1+N)^{a_k}\rho^r$.  Multiplication by the
uniformly bounded fixed derivatives of $G_N$ proves
\eqref{eq:v77-middle-rate}.

Equation~\eqref{eq:v64-green-reflection} and its fixed differentiated
versions give $|(G_N)_{ij}|\le C_kq_G^{|i-j|}$ after a strict
enlargement of the base to absorb derivative losses.  Opposite retained
ends are separated by at least $N-2r$.  Summing their matrix entries
costs at most a polynomial in $N$; multiply by the uniformly bounded
trace norm of the appropriate $\Delta_E$ block.  The inequality
$\|AB\|_2\le\|A\|\|B\|_2$ and $\|B\|_2\le\|B\|_1$
give \eqref{eq:v77-cross-rate}.  The same estimates within each end
give its uniform trace norm.

On the left block, the gluing estimate used for
\eqref{eq:v64-half-actions} bounds the difference between the finite
trajectory and the future half-line by the opposite-end weight
$C_k(1+N)^{a_k}\rho^{N-i}$, with a smaller gluing residual.
Summing on $i\le r$ gives $C_k(1+N)^{a_k}\rho^{N-r}$.
The reflected Green error on this block is at most
$C_k(1+N)^{a_k}q_G^{2(N-r)}$ by
\eqref{eq:v64-green-reflection}.  Multiplication by the localized
perturbation transfers both bounds to trace norm.  Extending the left
block to the entire half-line costs $C_k(1+N)^{a_k}\rho^r$:
perturbation columns have exponentially summable tails, and rows beyond
$r$ applied to retained columns are summed using $q_G<\rho$.
The right block is identical with the indices reversed.  This proves
\eqref{eq:v77-diagonal-rate} for the actual operators in
\eqref{eq:v64-half-amplitude}.

Choose the collar so that $\|G_N\|\,\|\Delta_N\|$ and the
corresponding half-line products have a common bound strictly below one.
Compression does not increase the operator norm of $G_N$ or of
$\Delta_N$.  Let $J_E=\diag(I,-I)$ on the two end spaces.  Then
\[
 \mathsf B_N+tE_N
 =\tfrac{1+t}{2}T_E+\tfrac{1-t}{2}J_ET_EJ_E,
 \qquad 0\le t\le1,
\]
so this whole segment has the same operator margin.  The deletion and
half-line comparison segments also have that margin by the preceding
product bounds and convexity.  Fixed parameter derivatives have bounded
trace norms; summation and product rules above give exactly the stated
fixed-order bounds, without an order-uniform assertion.

Apply Proposition~\ref{prop:v76-determinant}.  The deletion term is
\eqref{eq:v77-middle-rate}; the diagonal logarithm errors follow from
the trace-norm segment estimate and \eqref{eq:v77-diagonal-rate};
the off-diagonal error is the square of
\eqref{eq:v77-cross-rate}.  Enlarging a fixed polynomial exponent gives
\eqref{eq:v77-concrete-determinant}.  Finitely many small $N$ are
absorbed in the constants.  The local edge sum in
\eqref{eq:v64-relative-log} is controlled by the same gluing and tail
estimates, so the corollary applies before any probability normalization.
\end{proof}

The scalar factor in the subsequent conditional law is still
$D_{N,b}$ and cancels exactly, as in Corollary
\ref{cor:v76-relative-conditioning}.  This concrete instantiation is
separate from the finite-action reconstruction: the latter uses the
exact finite density rather than an $N\to\infty$ determinant limit.
The phase-weighted envelope also remains a refinement of the periodic
inverse, with its unweighted cost $\kappa(w)$; lowering its sufficient
differentiation order need not lower that cost.
